\begin{mistake}{2026-05-13}{04}

\begin{problemcontent}
曲线 \( y = \dfrac{1}{|x|-1} + \ln(1+e^{-x}) \) 的渐近线的条数为（\quad）

A. 1\qquad B. 2\qquad C. 3\qquad D. 4
\end{problemcontent}

\begin{solutioncontent}
渐近线需分三类逐一检查：铅直、水平、斜。

\textbf{1. 铅直渐近线}

\( |x|-1=0 \Rightarrow x=\pm1 \)，在这两点处 \( y\to\infty \).

铅直渐近线：\( x=1 \)，\( x=-1 \)（2条）。

\textbf{2. \( x\to+\infty \)（水平渐近线）}
\[
\frac{1}{x-1}\to 0,\quad \ln(1+e^{-x})\to 0 \implies y\to 0
\]
水平渐近线：\( y=0 \)（1条）。

\textbf{3. \( x\to-\infty \)（斜渐近线）}

此时 \( \ln(1+e^{-x})\approx -x\to+\infty \)，非水平，需检查斜渐近线。

斜率：
\[
k = \lim_{x\to-\infty}\frac{y}{x} = \lim_{x\to-\infty}\frac{\ln(1+e^{-x})}{x} \approx \frac{-x}{x} = -1
\]

截距：
\[
b = \lim_{x\to-\infty}\bigl[y-(-x)\bigr] = \lim_{x\to-\infty}\left[\frac{1}{-x-1}+\ln(1+e^{-x})+x\right]
\]
令 \( t=-x\to+\infty \)：
\[
= \lim_{t\to+\infty}\left[\frac{1}{t-1}+\ln(1+e^{t})-t\right] = 0 + 0 = 0
\]
其中 \( \ln(1+e^t)-t = \ln(1+e^{-t})\to 0 \)。

斜渐近线：\( y=-x \)（1条）。

共 \( 2+1+1=4 \) 条，选 D。
\end{solutioncontent}

\mistakeknowledge{
\begin{itemize}
  \item \textbf{核心公式/定理}：渐近线三类判定——铅直：\( \lim_{x\to a}y=\infty \)；水平：\( \lim_{x\to\infty}y=c \)；斜：\( k=\lim_{x\to\infty}\frac{y}{x}\neq0 \)，\( b=\lim_{x\to\infty}(y-kx) \) 存在。
  \item \textbf{易错点}：\( x\to-\infty \) 时 \( \ln(1+e^{-x})\approx -x \) 趋于无穷，不是水平渐近线，必须继续检查斜渐近线；\( |x| \) 导致铅直渐近线有对称两条。
  \item \textbf{关联知识}：同一方向上水平与斜渐近线互斥；左右两侧需分别讨论，可得不同类型渐近线。
\end{itemize}}

\end{mistake}
